\scanpage{16}
ces propriétés, on montre que, posant $M(h)=h(U)$, on \deleted{\unsure{est}} a
\[
 \alpha':\mathrm{Hom}(h,h')\xrightarrow{\sim}\mathrm{Hom}_{R^o}(\deleted{\unsure{h}}M(h),M(h'))
\]
où $R=\mathrm{Hom}_C(U,U)$, $R^o=$ anneau inverse de $R$.

Pour le voir, on définit un hom
\[
 \beta':\mathrm{Hom}_{R^o}(M(h),M(h'))\longrightarrow\mathrm{Hom}(h,h')
\]
Détails laissés au lecteur. Donc le foncteur $T':h\longrightarrow M(h)$ identifie la catégorie des foncteurs \deleted{Posons donc, pour \unsure{les}} envisagés à une sous-catégorie de $\unsure{\mathscr G|\mathfrak G|\mathcal Y}^{R^o}$ ($R$-mod à droite). On va définir un ``foncteur inverse'' $M\longrightarrow h_M$ de $\mathscr G^{R^o}$ dans la catégorie de foncteurs envisagée. On va définir directement le foncteur $T'':\deleted{\unsure{C}}\mathscr G^{R^o}\longrightarrow C$, ce sera simplement
\[
 T''(M)=M\otimes_R U
\]
\margin{\begin{tikzcd}[ampersand replacement=\&,row sep=small]
 C\arrow[d,"{T\ \text{``inj''}}"]\\
 \substack{\text{foncteurs contr.}\\ C\to\mathscr G\ \text{qui }\unsure{\text{tr}|\text{ch}}}\arrow[d,"{T'\ \text{``inj''}}"]\\
 \mathscr G^{R^o}\arrow[uu,bend left=65,"T''"]
\end{tikzcd}}
On va montrer que \deleted{si $h$ est un foncteur $C^o\longrightarrow\mathscr G$ du type envisagé, alors} \inserted{$M=T'T(T''M)$} donc $T'T$ est ``surjectif'', donc $T$ et $T'$ sont des ``équivalences''. On doit donc montrer que
\[
 \left\{\begin{aligned}
 &\mathrm{Hom}_C(U,M\otimes_R U)\simeq M\\
 &\text{si}\quad R=\mathrm{Hom}_C(U,U)\\
 &\phantom{\text{si}\quad} M\in\mathscr G^{R^o}
 \end{aligned}\right.
\]
ce qui est laissé au \unsure{soin|soins} du lecteur \ldots

\deleted{\unsure{L}} Ainsi, les foncteurs \deleted{\unsure{covariants}} \inserted{\unsure{contravariant|contravariants}} exacts à gauche $C\longrightarrow\mathscr G$, \deleted{transformant \unsure{une|les}} limites inductives en limites projectives sont essentiellement \inserted{des $h_B$}
% archive sheet number: 15
